% !TEX root = ../notes.tex

\documentclass[../notes.tex]{subfiles}

\begin{document}

\section{April 6}
Today, we write down some Vogan diagrams.

\subsection{Linear Algebra over the Quaternions}
To handle types the other classical types $B$, $C$, and $D$, we need to do a little quaternionic linear algebra. Recall the quaternions $\HH$ is a non-commutative quadratic extension of $\CC$ by an element $j$ for which $j^2=-1$ and $ij=-ji$.
\begin{notation}
	We define $\op{GL}_n(\HH)$ to be the group of $\RR$-linear automorphisms of $\HH^n$.
\end{notation}
\begin{remark}
	The group $\op{GL}_n(\HH)$ is a real form of $\op{GL}_{2n}(\CC)$.
\end{remark}
\begin{definition}[determinant]
	We define a determinant map $\det\colon\op{GL}_n(\HH)$ by sending $\alpha\in\op{GL}_n(\HH)$ to $\det\alpha\coloneqq\det_\CC(\alpha;\HH^n)$.
\end{definition}
\begin{remark} \label{rem:det-is-positive}
	The determinant turns out to output to $\RR^+$. It is enough to check this for diagonal matrices and unipotent upper-triangular matrices (by some row-reduction). There is nothing to say for unipotent matrices, and the case of diagonal matrices reduces to $n=1$. In this case, we are looking at the determinant map $\det\colon\HH^\times\to\CC^\times$, which can be computed directly. Alternatively, we can check that $\HH^\times\subseteq\op{GL}_2(\CC)$ turns out to be given by $\op{SU}(2)\times\RR^\times$, all of whose determinants go to $\RR^+$ because the diagonal $\RR^\times$ piece is squared.
\end{remark}
The moral of having a determinant is that we can define a special linear group.
\begin{notation}
	We define $\op{SL}_n(\HH)$ to be the subgroup of $\op{GL}_n(\HH)$ of elements with trivial determinant.
\end{notation}
\begin{remark}
	By \Cref{rem:det-is-positive}, there is a short exact sequence
	\[1\to\op{SL}_n(\HH)\to\op{GL}_n(\HH)\to\RR^+\to1.\]
	The same arguments over $\RR$ or $\CC$ tell us that we may take the differential to get a short exact sequence
	\[1\to\mf{sl}_n(\HH)\to\mf{gl}_n(\HH)\to\RR\to0\]
	of Lie algebras. One can check that the right-hand map sends $X$ to $\Re\tr X$; for example, this map should vanish on strictly upper-triangular or strictly lower-triangular matrices, so it reduces to a calculation on $n=1$, where it can be computed directly.
\end{remark}
We will also want to have some sesquilinear forms on $\HH$.
\begin{definition}[sesquilinear]
	Fix a vector space $V$ over $\HH$. Then a \textit{sesquilinear} form on $V$ is an $\mathbb R$-bilinear map $(-,-)\colon V\times V\to\HH$ for which
	\[(x\alpha,y\beta)=\ov\alpha(x,y)\beta\]
	for any $x,y\in V$ and $\alpha,\beta\in\HH$. We say that $(-,-)$ is \textit{Hermitian} if and only if $(x,y)=\overline{(y,x)}$, and it is \textit{skew-Hermitian} if and only if $(x,y)=-\overline{(y,x)}$.
\end{definition}
One can classify these forms as follows.
\begin{proposition}
	Fix a vector space $V$ over $\HH$.
	\begin{listalph}
		\item Every non-degenerate Hermitian form on $V$ is equivalent to
		\[(x,y)=(\overline x_1y_1+\cdots+\overline x_py_p)-(\overline x_{p+1}y{p+1}+\cdots+\overline x_{p+q}y_{p+q}),\]
		where $p+q=\dim_\HH V$.
		\item Every non-degenerate skew-Hermitian form is equivalent to
		\[(x,y)=\overline x_1jy_1+\cdots+\overline x_njy_n.\]
	\end{listalph}
\end{proposition}
\begin{proof}
	Omitted. This is similar to the argument for complex vector spaces.
\end{proof}
\begin{definition}[signature]
	Fix a non-degenerate Hermitian form $B$ on a quaternionic vector space $V$. Then its \textit{signature} is the pair $(p,q)$ where $B$ is equivalent to the standard form
	\[(x,y)=(\overline x_1y_1+\cdots+\overline x_py_p)-(\overline x_{p+1}y{p+1}+\cdots+\overline x_{p+q}y_{p+q}).\]
\end{definition}
\begin{remark} \label{rem:decompose-quaternionic-sesquilinear}
	Given a quaternionic sesquilinear form $B$ on $V$, then we can define $B_1,B_2\colon V\times V\to\CC$ so that $B=B_1+jB_2$, where each are $\RR$-bilinear. In fact, it turns out that $B_1$ is $\CC$-sesquilinear, and $B_2$ is bilinear, and they satisfy $B_2(x,y)=B_1(xj,y)$. Furthermore, $B$ is Hermitian of signature $(p,q)$ if and only if $B_1$ is Hermitian of signature $(2p,2q)$ and $B_2$ is skew-symmetric. Similarly, $B_1$ is skew-Hermitian if and only if $B_1$ is skew-Hermitian of signature $(n,n)$ and $B_2$ is symmetric.
\end{remark}
Our sesquilinear forms now have unitary groups.
\begin{definition}[unitary group]
	Fix a sesquilinear form $B$ of a quaternionic vector space $V$. Then we define the \textit{unitary group} $\op U(B)$ to be the group of linear transformations of $V$ which preserve $B$.
\end{definition}
\begin{example}
	Suppose that $B$ is Hermitian of signature $(p,q)$. Then \Cref{rem:decompose-quaternionic-sesquilinear} grants us a decomposition $B=B_1+jB_2$ where $B_1$ is Hermitian of signature $(p,q)$ and $B_2$ is skew-symmetry. It follows that
	\[\op U(B)=\op U(2p,2q)\cap\op{Sp}(2n).\]
	We will denote this group by $\op U(p,q,\HH)$ and its Lie algebra by $\mf{u}(p,q,\HH)$.
\end{example}
\begin{example}
	Suppose that $B$ is skew-Hermitian. Then \Cref{rem:decompose-quaternionic-sesquilinear} again allows us to decompose $B$, and we find that
	\[\op U(B)=\op U(n,n)\cap\op O(2n).\]
	We will denote this group by $\op O^*(2n)$ and its Lie algebra by $\mf{so}^*(2n)$.
\end{example}

\subsection{Real Forms of Classical Types}
We are now ready to write down all real forms in classical types. We already handled type $A$ in \Cref{exe:real-form-a}.
\begin{exe} \label{exe:real-form-b}
	We classify the real forms of type $B_n$.
\end{exe}
\begin{proof}
	Here, $\mf g=\mf{so}(2n+1)$. The Dynkin diagram looks like
	% https://q.uiver.app/#q=WzAsNSxbMCwwLCJcXGJ1bGxldCJdLFsxLDAsIlxcYnVsbGV0Il0sWzIsMCwiXFxjZG90cyJdLFszLDAsIlxcYnVsbGV0Il0sWzQsMCwiXFxidWxsZXQiXSxbMCwxLCIiLDAseyJzdHlsZSI6eyJoZWFkIjp7Im5hbWUiOiJub25lIn19fV0sWzEsMiwiIiwwLHsic3R5bGUiOnsiaGVhZCI6eyJuYW1lIjoibm9uZSJ9fX1dLFsyLDMsIiIsMCx7InN0eWxlIjp7ImhlYWQiOnsibmFtZSI6Im5vbmUifX19XSxbMyw0LCIiLDAseyJsZXZlbCI6Mn1dXQ==&macro_url=https%3A%2F%2Fraw.githubusercontent.com%2FdFoiler%2Fnotes%2Fmaster%2Fnir.tex
	\[\begin{tikzcd}[cramped]
		\bullet & \bullet & \cdots & \bullet & \bullet
		\arrow[no head, from=1-1, to=1-2]
		\arrow[no head, from=1-2, to=1-3]
		\arrow[no head, from=1-3, to=1-4]
		\arrow[Rightarrow, from=1-4, to=1-5]
	\end{tikzcd}\]
	which has no automorphisms, so all forms are in the same inner class. Now, we are hunting for $\theta^2=1$. Because $G_{\mathrm{ad}}=\op{SO}(2n+1)$, we find that $\theta$ acts by some matrix which can be taken to commute with our choice of maximal torus. Then $\theta$ lives in the corresponding torus of $\op{SO}(2n+1)$ (because it fixes the Dynkin diagram), so it diagonalizes appropriately. By using the Weyl group, we can rearrange the $(+1)$s and the $(-1)$s in its diagonalization, so we find that $\theta$ is given by conjugation by $1_p\oplus(-1)_q$ for some $p$ and $q$ with $p+q=n$. This gives the real form $\mf{so}(2p,2q+1)$, where $0\le p\le n$ and $q$ is chosen so that $p+q=n$. Finding $\theta$-stable elements reveals that $\mf k=\mf{so}(2p)\oplus\mf{so}(2q+1)$.
\end{proof}
\begin{exe}
	We classify the real forms of type $C_n$.
\end{exe}
\begin{proof}
	Once again, the Dynkin diagram has no automorphisms. Here, $G_{\mathrm{ad}}=\op{Sp}_{2n}(\CC)/\{\pm1\}$, so $\theta^2=1$ has two cases for a lift $\widetilde\theta$ to $\op{Sp}_{2n}(\CC)$. One can always change basis (as in \Cref{exe:real-form-b}) so that $\widetilde\theta$ is diagonal and has appropriately arranged diagonal elements.
	\begin{itemize}
		\item If $\widetilde\theta^2=1$, then we may assume $\widetilde\theta=1_p\oplus(-1_q)$. This gives the real form $\mf{u}(p,q,\HH)$, which has $\mf k^c=\mf u(p,\HH)\times\mf u(p,\HH)$ so that $\mf k=\mf{sp}_{2p}(\CC)\times\mf{sp}_{2p}(\CC)$.
		\item If $\widetilde\theta^2=-1$, then one can define $V=\CC^{2n}$ into two prices $V(i)$ and $V(-i)$, which one can see are both isotropic subspaces. Each must be at most $n$-dimensional (over $\CC$) because this is the largest size of an isotropic subspace, so for dimension reasons, they are both $n$-dimensional. Thus, $\mf k=\mf{gl}_n(\CC)$, namely fixing the eigenspaces. (The action on $V(-i)$ is determined by the action on $V(i)$ by definition of $\mf{sp}(2n)$.) Note that the real form is $\mf{sp}_{2n}(\RR)$ by exhaustion.
		\qedhere
	\end{itemize}
\end{proof}
\begin{exe}
	We classify the real forms of type $D_n$.
\end{exe}
\begin{proof}
	Let's start with the compact inner class. Here, $G_{\mathrm{ad}}=\op{SO}_{2n}(\CC)/\{\pm1\}$, so $\theta^2=1$ has two cases for a lift $\widetilde\theta$ to $\op{SO}_{2n}(\CC)$. As usual, we may assume that $\widetilde\theta$ is given by conjugation by a diagonal element.
	\begin{itemize}
		\item If $\widetilde\theta^2=1$, then we may assume that $\widetilde\theta=1_p\oplus(-1_q)$. This gives the real form $\mf{so}(2p,2q)$, which has $\mf k=\mf{so}_{2p}(\CC)\times\mf{so}_{2q}(\CC)$.
		\item If $\widetilde\theta^2=-1$, then we may assume that $\widetilde\theta=i1_p\oplus(-i1_q)$. But these eigenspaces subspaces are isotropic, so $p\le n$ and $q\le n$, so $p=q=n$ is forced by $p+q=n$. It follows that $\mf k=\mf{gl}_n$ by acting on one of the eigenspaces, and one finds that the real form is $\mf{so}^*(2n)$.
	\end{itemize}
	There is also another inner class corresponding to when $\theta$ is nontrivial.\footnote{This inner class is not necessarily the split inner class. Indeed, for $D_n$, the parity of $n$ determines if the compact inner class is split.} In this case, we see that $\theta$ actually lifts to conjugation by an element of $\op O(n)$ with determinant $-1$. Thus, $\theta$ can be taken to be a diagonal matrix $1_p\oplus(-1_q)$ where $p+q=n$ and $q$ is odd. Arguing as in the first point above, we see that we get the real form $\mf{so}(2p+1,2q-1)$, and $\mf k=\mf{so}(2p+1)\times\mf{so}(2q-1)$.
\end{proof}
\begin{remark}
	For $D_n$ at $n=3$, one can compare our calculation with $A_3$. By comparing $\mf k$s, we find $\mf{so}(6,\RR)=\mf{su}(4)$ because are the compact forms. By comparing $\mf k$s, we also find that $\mf{so}(4,2)=\mf{su}(2,2)$ and $\mf{so}(3,3)=\mf{sl}(4,\RR)$ and $\mf{so}(5,1)=\mf{sl}(2,\HH)$ and $\mf{so}^*(6)=\mf{su}(3,1)$.
\end{remark}

\end{document}